% InitechBase Volume Two: Programming with InitechBase, Index -- part
% 1040-0000-11 (PUB-REF-0001 Sec 12, DS row BAS-V2).  Family F-ATE; model
% the dBASE III PLUS 1.1 Volume Two binder, Programming section (annex Sec
% 1, 8.2 P, 10: one worked application built module by module; syntax
% questions sent back to Using).  The application is the Initech
% Accounting System on the Sample Programs and Utilities diskette; its
% month-end programs Y2KACCT.PRG and INTPOST.PRG carry the bodies of
% harness/diff/dbf_diff/canon/y2k_accounting.prg and salami.prg byte for
% byte (test-pubs-samir-programs), documented as the shipped programs.  Bead
% initech-ipu5.8.  Ground truth (Law 1): os/samir/** (samir_main.c the dot
% prompt and the message catalogue; cmd/*.c the commands; core/fn_builtins.c
% and core/eval.c the functions; fs/*.c the files), spec/samir/*.
% This volume teaches; it registers no entries (the reference is Using
% InitechBase), and every inline command, function and message it names is
% checked against the shipped surface (test-pubs-commands); every example is
% replayed through the shipped dot
% prompt against a fresh copy of the Sample Programs and Utilities diskette
% (pubs/examples/samir/disk, harness/pubs/samir_replay.c, make
% test-pubs-examples).  No version is printed: the shipped build displays
% none (D-3, PN-03).
\documentclass[ashtontate,binder]{initech-manual}

\ProductName{InitechBase}
\DocTitle{Volume Two}
\PartNumber{1040-0000-11}
\FamilyName{Information Management Family}
\FamilyColor{FamilyInformation}
\Edition{First Edition (March 1992)}
\EditionApplies{This edition applies to the release of InitechBase that it
  accompanies and to all subsequent releases until otherwise indicated in new
  editions.}
\CopyrightYears{1991, 1992}
\Revision{R01}
\PrintCode{A1-3/92}
\Trademarks{InitechOS, InitechDOS, FLAIR, InitechBase, InitechCalc and
  InitechPaint are trademarks of Initech Systems Corporation.  Turbo Initech
  is a registered trademark of Initech Systems Corporation.  Intel, 80386 and
  80486 are trademarks of Intel Corporation.  All other product names are
  trademarks of their respective holders.}
% Operator directive 2026-10-17 (the design lane's class macro wins once it
% merges): the corporate motto at the placements the directive names.
\providecommand{\IsThisGoodForTheCompany}{\par\medskip\noindent{\sffamily Is This Good for the \textbf{\textit{COMPANY}}?}\par\noindent\rule{\linewidth}{0.5pt}\medskip}

% A defined term at its first use (TY-02, NT-34): bold, indexed, and in the
% Glossary (CP-03); the optional argument is the singular index form.
\NewDocumentCommand\term{om}{\textbf{#2}\IfValueTF{#1}{\index{#1}}{\index{#2}}}
% What the reader types (NT-15), breakable at its blanks: the dot-prompt
% commands of the lessons are longer than a DOS command.
\renewcommand\type[1]{\texttt{\bfseries #1}}
\newcommand\bs{\textbackslash}

\XbaseIndexMentions
\MasterIndexPartner{initechbase-volume-one}
% AN-26 / D-15: the colophon names the replay route of the examples.
\newcommand\ColophonInitechBase{}
\begin{document}
\FrontMatterFolios
\MakeTitlePage
% Front matter of Volume Two: edition notice (PUB-REF-0001 AN-20, fixed
% order; boilerplate newly written, VO-12); the Introduction of the
% section, with its conventions, is p0-introduction.tex (AN-21 F-ATE).
% The motto follows the edition notice (operator directive 2026-10-17).
\Gates{test-pubs-build}
\begin{TitleVerso}
{\imLegalFont
\textbf{The following paragraph does not apply to any country where such
provisions are inconsistent with local law:} INITECH SYSTEMS CORPORATION
PROVIDES THIS PUBLICATION ``AS IS'' WITHOUT WARRANTY OF ANY KIND, EITHER
EXPRESS OR IMPLIED, INCLUDING, BUT NOT LIMITED TO, THE IMPLIED WARRANTIES OF
MERCHANTABILITY AND FITNESS FOR A PARTICULAR PURPOSE.  Some states do not
allow the exclusion of implied warranties, so the above exclusion may not
apply to you.\par
This publication may contain technical inaccuracies or typographical errors.
The information in it is revised from time to time, and the revisions are
published in new editions.  Initech Systems Corporation may change the
products and programs described in this publication at any time without
notice.\par
Requests for copies of Initech publications and for technical information
about Initech products should be made to your authorized Initech dealer or
your Initech representative.\par}

{\footnotesize
Reader's Comment Forms are provided at the back of this publication.  If the
forms have been removed, address comments to: \InitechAddressLine.
\draft{PUB-REF-0001 O-5 and D-12: the address and permit number are to be
checked against real businesses before print}
Initech may use or distribute the information you supply in any way it
believes appropriate without incurring any obligation to you.\par
The software described in this publication is furnished under the Initech
Program License Agreement and may be used and copied only as that agreement
permits.  The sample databases, report forms and programs on the Sample
Programs and Utilities diskette may be copied and changed for use with
InitechBase.\par}
\end{TitleVerso}


\tableofcontents

\BeginSection{P}{Programming with InitechBase}
% Programming with InitechBase, Introduction (PUB-REF-0001 AN-21 F-ATE:
% what the section is, who should read it, chapter-by-chapter list,
% Conventions Used in This Manual; annex Sec 1, 8.1).  Folios PI-n.
\SectionIntroduction{Introduction}
\Gates{test-pubs-build}
\doctitle{Programming with InitechBase} teaches you to write programs: files of
InitechBase commands that do a task from beginning to end whenever you run
them.  It is written for readers who have worked through
\doctitle{Learning InitechBase} and know the commands at the dot prompt.  You
need no experience of programming.

The section builds one application, the Initech Accounting System, a chapter
at a time.  The system keeps the accounts receivable ledger of Initech
Systems Corporation in the database \file{INVOICE.DBF}: one record for each
invoice, with its number, the customer, the amount, the due date and whether
it has been paid.  Its programs list and find invoices, enter new ones, print
the invoice register, age the open invoices at the end of the month, and post
the monthly finance charges.  Every program is on the Sample Programs and
Utilities diskette, and every program is printed in full in this section.

For the exact form of a command or function, this section refers you to
\doctitle{Using InitechBase} in Volume One.

\section{How This Section Is Organized}
\SideHead{Chapter 1}Programming in a Nutshell: what a program is, and the first
program of the system.

\SideHead{Chapter 2}Program Structure and Flow: sequence, decisions and loops.

\SideHead{Chapter 3}Memory Variables: totals, types and the life of a memory
variable.

\SideHead{Chapter 4}Procedures and Parameters: dividing a program into parts.

\SideHead{Chapter 5}Functions for Fields and Memory Variables: building and
formatting values.

\SideHead{Chapter 6}Communicating with the User: questions and answers.

\SideHead{Chapter 7}Working with the Database: finding, adding and changing
records in a program.

\SideHead{Chapter 8}Printing: the invoice register.

\SideHead{Chapter 9}The Aging Report: the month-end aging of the open invoices.

\SideHead{Chapter 10}Finance-Charge Posting: the monthly finance charges and
the rounding adjustment.

\SideHead{Chapter 11}The Main Menu: the program the clerk uses every day.

\SideHead{Chapter 12}Handling Errors: what a program does when a command
fails.

\SideHead{Chapter 13}Testing and Debugging: finding the mistake in a program.

\SideHead{Chapter 14}Working with Two Databases: a lookup in a second work
area.

\SideHead{Chapter 15}Reports from Programs: pages, headings and page numbers.

\SideHead{Chapter 16}Working with Dates: years of service.

\SideHead{Chapter 17}Working with Characters: names for letters.

\SideHead{Chapter 18}Changing Many Records: raises by department.

\SideHead{Chapter 19}Memo Fields in Programs: the remarks of the roster.

\SideHead{Chapter 20}Order and Filters in Programs: two orders and a filter.

\SideHead{Chapter 21}Searching: a search loop with \cmd{LOCATE}.

\SideHead{Chapter 22}Putting It Together: the Initech Accounting System as a
whole, and the month-end procedure.

\section{Conventions Used in This Manual}
\SideHead{Programs}Program listings are printed in bold on a shaded panel,
line for line as they are stored on the sample diskette.

\SideHead{Examples}The panels of examples show the screen: the dot prompt, a
period and a blank, the lines you type in bold, and what InitechBase displays.
Every example was run against a fresh copy of the sample diskette; your
results agree with the book if your files are as they were on the diskette.

\SideHead{Names}The programs of this section begin the names of their memory
variables with M, so that a memory variable is never taken for a field.

\SideHead{Dates}The examples were run on the 1st of March 1992, the date
\fn{DATE()} returns in them.  The accounting programs take their reporting
date from a memory variable instead, so that a month can be closed again
later.

% Programming with InitechBase, Chapter 1: Programming in a Nutshell
% (AN-16; annex Sec 8.2 P1).  Panels replayed: examples/p1-*.dot (dot-prompt
% route) and ../../examples/samir/salary (the program route of
% harness/pubs/samir_replay.c, samir_do).  Sources samir_main.c
% repl_try_do_file, cmd/proc.c, cmd/flow.c.
\chapter{Programming in a Nutshell}
\Gates{test-pubs-examples, test-interp-proc, test-interp-flow}
A program is a list of commands written in a file.  When you run the program,
InitechBase carries out the commands one after another, as if you were typing
them at the dot prompt, but faster and without mistakes.  In this chapter you
see what a program is made of, run the first program of the Initech Accounting
System, and read it line by line.

\begin{ChapterCovers}
\item From the dot prompt to a program
\item Program files
\item Running a program with \cmd{DO}
\item Reading the first program of the system
\end{ChapterCovers}

\ATEHead{Preparing for This Chapter}
Copy the Sample Programs and Utilities diskette to the disk you work on, as
\doctitle{Getting Started} describes, and start InitechBase from it.  Copy the
diskette again before each chapter: the programs of later chapters change the
ledger.

\ATEHead{From the Dot Prompt to a Program}
At the dot prompt you can look at the ledger an invoice at a time:
\begin{Steps}
\item Type \type{USE INVOICE}, then \type{GO TOP}.
\item Type \type{? INVNO, CUST, AMOUNT}, then \type{SKIP}, then \type{? INVNO, CUST,
AMOUNT}.
\end{Steps}
\DotExample{examples/p1-dot}
To look at every open invoice this way you would type the same two commands
for each record, and decide each time whether the invoice is paid.  A program
does the typing and the deciding for you.

\ATEHead{Program Files}
A program file is a text file with the extension \file{.PRG}, with one
InitechBase command on each line.  You write it with any editor that saves
plain text, such as the InitechDOS line editor \file{EDLIN}.  The commands are
those you use at the dot prompt, with a few more that exist only in programs:
\cmd{IF}, \cmd{DO CASE} and \cmd{DO WHILE}, which decide and repeat, and
\cmd{PROCEDURE} and \cmd{RETURN}, which divide a program into parts.

A program uses the databases that are open when it runs.  Open them at the dot
prompt with \cmd{USE} before you run the program: \cmd{USE}, \cmd{CLOSE} and
\cmd{QUIT} are dot prompt commands, and a program does not use them.

Here is a very small program.  It sets the decimal places, stores a salary in
a memory variable and displays it with an allowance added:
\VerbatimInput[fontsize=\small,xleftmargin=4mm]{../../examples/samir/salary.prg}
When it runs, it displays:
\begin{Transcript}{../../examples/samir/salary}
\end{Transcript}

\ATEHead{Running a Program}
\cmd{DO} runs a program.  You need not type the extension.  The first program of
the Initech Accounting System, \file{OPENINV.PRG}, lists the invoices that have
not been paid:
\begin{Steps}[resume]
\item Type \type{USE INVOICE}.
\item Type \type{DO OPENINV}.
\ATEResponse{The program lists the open invoices and counts them.}
\end{Steps}
\DotExample{examples/p1-openinv}
When the program ends, InitechBase shows the dot prompt again.

\ATEHead{Reading the Program}
\ProgramListing{OPENINV.PRG}
\SideHead{Comments}The first lines begin with an asterisk.  They are comments:
InitechBase passes over them.  Every program of the system begins with its
name, its purpose and the database it expects.

\SideHead{Headings}Two \cmd{?} commands print the title and the column
headings.

\SideHead{The count}\texttt{MCOUNT = 0} creates the memory variable that counts
the open invoices.

\SideHead{The loop}\cmd{GO TOP} moves to the first record.  The commands
between \cmd{DO WHILE} and \cmd{ENDDO} are repeated as long as the end of the
file has not been reached; \cmd{SKIP}, at the end of the loop, moves to the
next record each time.

\SideHead{The decision}Inside the loop, \cmd{IF .NOT. PAID} chooses the open
invoices.  For each, \cmd{?} prints the invoice number, the customer and the
amount, and the count grows by one.  \fn{STR()} turns the amount into
characters, so that it can be joined to the other values with \texttt{+}.

\SideHead{The end}After the loop, the program prints the count and
\cmd{RETURN} ends it.

\ATEHead{Where Do You Go From Here?}
Chapter~2 looks more closely at the structures that decide and repeat.

\ATEHead{Summary}
\begin{enumerate}
\item A program is a text file of commands, with the extension \file{.PRG}.
\item Open the databases with \cmd{USE}, then run the program with \cmd{DO}.
\item Comment lines begin with an asterisk.
\item \cmd{DO WHILE .NOT. EOF()} \ldots{} \cmd{SKIP} \ldots{} \cmd{ENDDO} visits
  every record; \cmd{IF} chooses records.
\item \cmd{RETURN} ends the program.
\end{enumerate}

% Programming with InitechBase, Chapter 2: Program Structure and Flow
% (AN-16; annex Sec 8.2 P2).  Panels replayed (examples/p2-*.dot); sources
% cmd/flow.c (IF/ELSE/ENDIF, DO WHILE/ENDDO/LOOP/EXIT, DO CASE; guards must
% be Logical, #37).
\chapter{Program Structure and Flow}
\Gates{test-pubs-examples, test-interp-flow}
A program runs its commands in order, from the first line to the last, unless a
structure tells it otherwise.  InitechBase has three: \cmd{IF} chooses whether
to carry out a group of commands, \cmd{DO CASE} chooses one group among
several, and \cmd{DO WHILE} repeats a group.  This chapter describes each, with
the programs of the system that use them.

\begin{ChapterCovers}
\item Sequence
\item \cmd{IF}, \cmd{ELSE} and \cmd{ENDIF}
\item \cmd{DO CASE}
\item \cmd{DO WHILE}, \cmd{LOOP} and \cmd{EXIT}
\item Nesting and indentation
\item Conditions
\end{ChapterCovers}

\ATEHead{Preparing for This Chapter}
Start InitechBase with a fresh copy of the sample diskette, and open the ledger
with \type{USE INVOICE}.

\ATEHead{Sequence}
Without a structure, the commands of a program run one after the other, once
each.  The first five commands of \file{OPENINV.PRG} (Chapter~1) are a sequence:
two headings, a memory variable and a move to the first record.

\ATEHead{Decisions: IF}
\cmd{IF} is followed by a condition.  When the condition is true, the commands
up to \cmd{ELSE}, or up to \cmd{ENDIF} if there is no \cmd{ELSE}, are carried
out.  When it is false, the commands after \cmd{ELSE} are carried out instead.
Either way, the program goes on after \cmd{ENDIF}.

\SyntaxLine{\kw{IF} \xp{condition}}
\SyntaxLine{\ \ \xp{commands when true}}
\SyntaxLine{\kw{ELSE}}
\SyntaxLine{\ \ \xp{commands when false}}
\SyntaxLine{\kw{ENDIF}}

\file{INVTOTAL.PRG} adds every amount to the total of the ledger, and the paid
amounts to a second total:
\ProgramListing{INVTOTAL.PRG}

\ATEHead{Conditions}
A condition must be a logical value: a logical field such as \texttt{PAID}, a
comparison, a logical function such as \fn{EOF()} or \fn{FOUND()}, or these
joined with \texttt{.AND.}, \texttt{.OR.} and \texttt{.NOT.}.  InitechBase does
not treat a number or a character string as true or false.  A program whose
condition is of another type stops with
\xmsg{37}{Not a Logical expression.}

\ATEHead{Choices: DO CASE}
\cmd{DO CASE} holds several \cmd{CASE} lines, each with a condition.  The
commands of the first \cmd{CASE} whose condition is true are carried out, and
the program goes on after \cmd{ENDCASE}.  If no condition is true, the commands
after \cmd{OTHERWISE} are carried out.

\SyntaxLine{\kw{DO CASE}}
\SyntaxLine{\kw{CASE} \xp{condition 1}}
\SyntaxLine{\ \ \xp{commands}}
\SyntaxLine{\kw{CASE} \xp{condition 2}}
\SyntaxLine{\ \ \xp{commands}}
\SyntaxLine{\kw{OTHERWISE}}
\SyntaxLine{\ \ \xp{commands}}
\SyntaxLine{\kw{ENDCASE}}

The main menu of the system (Chapter~11) carries out the task the clerk
chooses with a \cmd{DO CASE}, one \cmd{CASE} for each choice.

\ATEHead{Counting by Class}
\file{DUECLASS.PRG} counts the invoices of the ledger in three classes of
amount.  The \cmd{DO CASE} inside the loop chooses the count for each record:
\ProgramListing{DUECLASS.PRG}
\DotExample{examples/p2-dueclass}
The order of the \cmd{CASE} lines matters.  An amount under 100 is also under
1000, but only the first \cmd{CASE} that is true is carried out, so it is
counted once.

\ATEHead{Loops: DO WHILE}
\cmd{DO WHILE} tests its condition; if it is true, the commands up to
\cmd{ENDDO} are carried out and the condition is tested again.  The loop ends
when the condition is false.

\SyntaxLine{\kw{DO WHILE} \xp{condition}}
\SyntaxLine{\ \ \xp{commands}}
\SyntaxLine{\kw{ENDDO}}

The loop of nearly every program of the system processes the records of the
ledger:

\SyntaxLine{\kw{GO TOP}}
\SyntaxLine{\kw{DO WHILE .NOT. EOF()}}
\SyntaxLine{\ \ \xp{commands for one record}}
\SyntaxLine{\ \ \kw{SKIP}}
\SyntaxLine{\kw{ENDDO}}

\cmd{LOOP} goes back to the test at once, leaving out the rest of the loop;
\cmd{EXIT} leaves the loop at once.  The menu of Chapter~11 is a loop that
never ends by its condition, \cmd{DO WHILE .T.}: \cmd{EXIT} ends it when the
clerk chooses to leave.

The memory variables of the program remain after it ends, so you can look at
them at the dot prompt:
\DotExample{examples/p2-openinv}

\ATEHead{Leaving a Loop Early}
\file{FIRSTBIG.PRG} looks for the first open invoice of 1000 or more.  It
passes over paid invoices with \cmd{LOOP} -- note the \cmd{SKIP} before it, or
the loop would test the same record for ever -- and leaves the loop with
\cmd{EXIT} as soon as it finds one:
\ProgramListing{FIRSTBIG.PRG}
\DotExample{examples/p2-firstbig}
The logical memory variable \texttt{MFOUND} remembers why the loop ended: at
the end of the file, or at an invoice found.

\ATEHead{Nesting and Indentation}
A structure may hold other structures: an \cmd{IF} inside a \cmd{DO WHILE}, as
in \file{OPENINV.PRG}, or a \cmd{DO CASE} inside an \cmd{IF}.  Each must be
closed inside the one that holds it.  InitechBase ignores blanks at the
beginning of a line, so indent the commands inside each structure by two
blanks: the programs of the system are written that way, and the structure is
then plain to see.

\begin{Tip}
Write the closing line of a structure -- \cmd{ENDIF}, \cmd{ENDCASE},
\cmd{ENDDO} -- as soon as you write its first line, then fill in the commands
between.
\end{Tip}

\ATEHead{Where Do You Go From Here?}
The programs of this chapter keep their totals in memory variables.  Chapter~3
tells you how memory variables are created and how long they last.

\ATEHead{Summary}
\begin{enumerate}
\item Commands run in sequence unless a structure decides or repeats.
\item \cmd{IF} \ldots{} \cmd{ELSE} \ldots{} \cmd{ENDIF} chooses between two groups.
\item \cmd{DO CASE} \ldots{} \cmd{ENDCASE} chooses the first group whose condition
  is true.
\item \cmd{DO WHILE} \ldots{} \cmd{ENDDO} repeats; \cmd{LOOP} and \cmd{EXIT} jump.
\item Every condition is a logical value.
\end{enumerate}

% Programming with InitechBase, Chapter 3: Memory Variables (AN-16; annex
% Sec 8.2 P3).  Panels replayed (examples/p3-*.dot); sources cmd/flow.c
% (STORE, =, the memvar table, 256 variables, names of 10), cmd/proc.c
% (PUBLIC, PRIVATE, scope), core/fn_builtins.c TYPE.
\chapter{Memory Variables}
\Gates{test-pubs-examples, test-interp-flow, test-interp-proc}
A program needs places to keep values that are not fields of a database: a
total, a count, an answer, a date.  These are memory variables.  This chapter
describes how they are created, what types they have, and how long they live.

\begin{ChapterCovers}
\item Creating memory variables
\item Types
\item Totals and counts
\item Public and private memory variables
\item Releasing memory variables
\end{ChapterCovers}

\ATEHead{Preparing for This Chapter}
Start InitechBase with a fresh copy of the sample diskette.

\ATEHead{Creating Memory Variables}
A memory variable is created when a value is first stored in it, by
\cmd{STORE} or by the equals sign:

\SyntaxLine{\kw{STORE} \xp{exp} \kw{TO} \xp{name list}}
\SyntaxLine{\xp{name} \kw{=} \xp{exp}}

\cmd{STORE} can give several memory variables the same value at once; the
programs of the system use it to set their totals to zero.  A name has up to 10
letters, digits and underscores and begins with a letter.  The system begins
every name with M: \texttt{MTOTAL}, \texttt{MCOUNT}, \texttt{MINVNO}.  A name that
is also the name of a field would always stand for the field.

\ATEHead{Types}
A memory variable takes the type of the value stored in it: character,
numeric, date or logical.  The function \fn{TYPE()} tells you the type of a
value:
\DotExample{examples/p3-types}
Storing a value of another type in the memory variable changes its type.

\ATEHead{Totals and Counts}
A total is a memory variable that starts at zero and grows by an amount for
each record; a count grows by one.  \file{INVTOTAL.PRG}, listed in Chapter~2,
keeps two totals:
\DotExample{examples/p3-invtotal}

\ATEHead{Public and Private Memory Variables}
The memory variables created by the main program of a file remain after it
ends, as if you had created them at the dot prompt.  Those created inside a
procedure (Chapter~4) are private to it: they are released when the procedure
returns.

\cmd{PUBLIC} declares memory variables that every part of every program can use
and that remain until they are released.  \file{INVTOTAL.PRG} declares its two
totals \cmd{PUBLIC}, so that the month-end journal entry can be made from them
after the program has run.  Declare a memory variable \cmd{PUBLIC} before you
store a value in it.

\cmd{PRIVATE} declares memory variables that belong to the procedure that
declares them.  A memory variable of the same name created by the caller is
hidden while the procedure runs and is unchanged afterwards.  A procedure that
uses a name its caller may also use, such as \texttt{MCOUNT}, declares it
\cmd{PRIVATE}.

\ATEHead{Releasing Memory Variables}
\cmd{RELEASE} removes memory variables, and \cmd{RELEASE ALL} removes all of them.
InitechBase holds up to 256 memory variables; a program that creates many in a
loop releases those it no longer needs.

\ATEHead{Where Do You Go From Here?}
Chapter~4 divides a program into procedures and passes values to them.

\ATEHead{Summary}
\begin{enumerate}
\item \cmd{STORE} and \texttt{=} create memory variables and change their values.
\item A memory variable takes the type of its value; \fn{TYPE()} tells it.
\item A total starts at zero and grows for each record.
\item \cmd{PUBLIC} memory variables remain; \cmd{PRIVATE} ones belong to their
  procedure.
\item \cmd{RELEASE} removes memory variables.
\end{enumerate}

% Programming with InitechBase, Chapter 4: Procedures and Parameters
% (AN-16; annex Sec 8.2 P4).  Panels replayed (examples/p4-*.dot); sources
% cmd/proc.c (PROCEDURE, DO ... WITH by value, PARAMETERS, RETURN, PUBLIC),
% samir_main.c (DO of a disk program only at the dot prompt).
\chapter{Procedures and Parameters}
\Gates{test-pubs-examples, test-interp-proc}
A calculation that a program needs more than once, or that makes the program
hard to read, is written once as a procedure and called by name.  This chapter
shows you how to write a procedure, pass values to it and get its result back.

\begin{ChapterCovers}
\item Procedures in a program file
\item Calling a procedure with \cmd{DO}
\item Passing values with \kw{WITH} and \cmd{PARAMETERS}
\item Returning a result
\end{ChapterCovers}

\ATEHead{Preparing for This Chapter}
Start InitechBase with a fresh copy of the sample diskette.

\ATEHead{Procedures}
The procedures of a program follow its main program in the same file.  Each
begins with a \cmd{PROCEDURE} line that names it and ends with \cmd{RETURN}.
The main program ends at the first \cmd{PROCEDURE} line, so a procedure runs only
when \cmd{DO} calls it.

A program calls the procedures of its own file.  A program file is run from the
dot prompt; one program does not run another program file.  The main menu of
the system (Chapter~11) therefore holds its procedures in its own file.

\ATEHead{Passing Values}
\kw{WITH} after the name of the procedure gives the values to pass, separated by
commas.  The first command of the procedure, \cmd{PARAMETERS}, names the memory
variables that receive them, in the same order.  The procedure receives copies:
what it does with its parameters does not change the values in the caller.

\file{LATEFEE.PRG} computes the late fee of an invoice: one per cent of the
amount for each full thirty days the invoice is past due.  The procedure
\texttt{FEE} receives the amount and the number of days:
\ProgramListing{LATEFEE.PRG}

\ATEHead{Returning a Result}
A procedure gives back its result in a memory variable declared \cmd{PUBLIC},
which remains when the procedure returns.  \texttt{FEE} stores the fee in
\texttt{MFEE}:
\begin{Steps}
\item Type \type{USE INVOICE}, then \type{GO 2}.
\item Type \type{DO LATEFEE}, then \type{? MFEE}.
\end{Steps}
\DotExample{examples/p4-latefee}
The parameters \texttt{PAMOUNT} and \texttt{PDAYS} are private to the procedure
and are released when it returns.

\begin{AteNote}
Pass character values to a procedure through memory variables or expressions
(\doctitle{Using InitechBase}, Chapter~5, \cmd{DO}).
\end{AteNote}

\ATEHead{Where Do You Go From Here?}
Chapter~5 looks at the functions that build the values programs print.

\ATEHead{Summary}
\begin{enumerate}
\item A procedure begins with \cmd{PROCEDURE} \xp{name} and ends with
  \cmd{RETURN}; it follows the main program in the file.
\item \cmd{DO} \xp{procedure} \kw{WITH} \xp{values} calls it;
  \cmd{PARAMETERS} \xp{names} receives the values.
\item A result is returned in a \cmd{PUBLIC} memory variable.
\end{enumerate}

% Programming with InitechBase, Chapter 5: Functions for Fields and Memory
% Variables (AN-16; annex Sec 8.2 P5).  Panels replayed (examples/p5-*.dot);
% sources core/fn_builtins.c (STR, LTRIM, TRANSFORM, DTOC, CTOD, CMONTH,
% IIF in core/eval.c).
\chapter{Functions for Fields and Memory Variables}
\Gates{test-pubs-examples, test-xbase-fn-a, test-xbase-fn-c, test-xbase-transform}
A program prints values in columns, joins them into lines, and converts them
from one type to another.  Functions do this work.  This chapter shows the
functions the Initech Accounting System uses to build its lines; Chapter~6 of
\doctitle{Using InitechBase} describes every function.

\begin{ChapterCovers}
\item Joining values into a line
\item Numbers as characters: \fn{STR()}, \fn{LTRIM()}, \fn{TRANSFORM()}
\item Dates as characters: \fn{DTOC()} and \fn{CTOD()}
\item Choosing a value: \fn{IIF()}
\item A formatted register
\end{ChapterCovers}

\ATEHead{Preparing for This Chapter}
Start InitechBase with a fresh copy of the sample diskette, and type
\type{USE INVOICE}.

\ATEHead{Joining Values}
\texttt{+} joins character values into one, which \cmd{?} then prints as a single
line.  Every part must be a character value: InitechBase does not join a number
or a date to characters.  Fields of type Character are joined as they are,
blanks and all, which keeps the columns of a listing in line.

\ATEHead{Numbers as Characters}
\fn{STR()} converts a number to characters of a fixed width, with a fixed number
of decimal places, right-aligned; it is the function that lines numbers up in
a column.  \fn{LTRIM()} removes the leading blanks when the number is to follow
a word.  \fn{TRANSFORM()} converts a number by a picture, in which each 9 holds a
digit:
\DotExample{examples/p5-numbers}

\ATEHead{Dates as Characters}
\fn{DTOC()} converts a date to characters in the form set by SET DATE;
\fn{CTOD()} converts characters to a date.  A date plus a number of days is a
date; \fn{CMONTH()} gives the name of its month:
\DotExample{examples/p5-dates}
The year of a date has two digits unless SET CENTURY is ON.  Two digits typed
into \fn{CTOD()} are a year of the nineteen hundreds; this is the date
convention of every program of the system.

\ATEHead{Choosing a Value}
\fn{IIF()} returns its second argument when its first is true, and its third when
it is false.  It chooses a word for a logical field without an \cmd{IF}.

\ATEHead{A Formatted Register}
\file{INVFMT.PRG} prints the ledger in fixed columns, with the status of each
invoice in words:
\ProgramListing{INVFMT.PRG}
\DotExample{examples/p5-invfmt}
Each line is built from the fields with \texttt{+}, in two steps to keep the
program lines short: the memory variable \texttt{MLINE} holds the first part.
\fn{DTOC()} converts the due date, \fn{TRANSFORM()} the amount, and \fn{IIF()}
turns \texttt{PAID} into a word, kept in the memory variable \texttt{MSTAT}.

\ATEHead{Where Do You Go From Here?}
Chapter~6 lets the program ask the user for what it needs.

\ATEHead{Summary}
\begin{enumerate}
\item \texttt{+} joins character values; numbers and dates are converted first.
\item \fn{STR()}, \fn{LTRIM()} and \fn{TRANSFORM()} turn numbers into characters.
\item \fn{DTOC()} and \fn{CTOD()} convert between dates and characters.
\item \fn{IIF()} chooses between two values.
\end{enumerate}

% Programming with InitechBase, Chapter 6: Communicating with the User
% (AN-16; annex Sec 8.2 P6).  Panels replayed (examples/p6-*.dot); sources
% cmd/proc.c (ACCEPT, INPUT, WAIT), cmd/query.c (? and ??).
\chapter{Communicating with the User}
\Gates{test-pubs-examples, test-interp-proc}
A program that does the same thing every time needs no questions.  Most need to
know something first: which invoice, which customer, which choice of the menu.
This chapter shows how a program asks, and how it answers.

\begin{ChapterCovers}
\item Displaying results: \cmd{?} and \cmd{??}
\item Asking for characters: \cmd{ACCEPT}
\item Asking for a value: \cmd{INPUT}
\item Asking for a key: \cmd{WAIT}
\item Finding an invoice by its number
\end{ChapterCovers}

\ATEHead{Preparing for This Chapter}
Start InitechBase with a fresh copy of the sample diskette, and type
\type{USE INVOICE INDEX INVNO}.

\ATEHead{Displaying Results}
\cmd{?} begins a new line and displays its values; \cmd{??} displays them where
the cursor is.  A program builds a line from several \cmd{??} commands, or from
one \cmd{?} and the \texttt{+} operator.  \cmd{?} with nothing after it prints an
empty line, which a program uses to separate its output.

\ATEHead{Asking for Characters}
\cmd{ACCEPT} displays a prompt and stores the line the user types in a memory
variable, as characters:

\SyntaxLine{\kw{ACCEPT} \xp{prompt} \kw{TO} \xp{memvar}}

The user need not type quotation marks.  Convert the answer to capitals with
\fn{UPPER()} before you compare it with fields kept in capitals.

\ATEHead{Asking for a Value}
\cmd{INPUT} evaluates the line the user types and stores its value, of whatever
type: a number for an amount, \texttt{.T.} or \texttt{.F.} for a yes or no.  A
character string must then be typed in quotation marks, so the system uses
\cmd{ACCEPT} for characters and \cmd{INPUT} for amounts.

\ATEHead{Asking for a Key}
\cmd{WAIT} takes one key, without \key{Enter}.  The menu of Chapter~11 uses it
for the clerk's choice.  \cmd{WAIT} without \kw{TO} only waits, to hold output
on the screen.

\ATEHead{Finding an Invoice}
\file{FINDINV.PRG} asks for an invoice number and displays the invoice:
\ProgramListing{FINDINV.PRG}
\begin{Steps}
\item Type \type{DO FINDINV}.
\ATEResponse{The program asks \texttt{Invoice number:}.}
\item Type \type{a1002} and press \key{Enter}.
\item Type \type{DO FINDINV} again and answer \type{A9999}.
\end{Steps}
\DotExample{examples/p6-findinv}
The answer is converted to capitals and looked up in the index \file{INVNO.NDX}
with \cmd{SEEK}.  \fn{FOUND()} decides which of the two messages the program
prints.

\ATEHead{Checking an Answer}
A program cannot trust what is typed.  \file{AMTCHECK.PRG} asks for an amount
until the answer is a number in the range it accepts:
\ProgramListing{AMTCHECK.PRG}
\DotExample{examples/p6-amtcheck}
\fn{TYPE()} tells the program whether the answer is a number before it is
compared; a character string typed in quotation marks is refused like a number
out of range.  The loop ends only when \texttt{MOK} is true.

\ATEHead{Where Do You Go From Here?}
Chapter~7 adds records to the ledger from the answers to questions.

\ATEHead{Summary}
\begin{enumerate}
\item \cmd{?} and \cmd{??} display; \cmd{?} alone prints an empty line.
\item \cmd{ACCEPT} asks for characters; \cmd{INPUT} for a value of any type;
  \cmd{WAIT} for one key.
\item Convert character answers with \fn{UPPER()} before comparing them.
\end{enumerate}

% Programming with InitechBase, Chapter 7: Working with the Database (AN-16;
% annex Sec 8.2 P10/P11).  Panels replayed (examples/p7-*.dot); sources
% cmd/query.c (SEEK, FOUND), cmd/mutate.c (APPEND BLANK, REPLACE, index
% maintenance), samir_main.c repl_do_use (INDEX).
\chapter{Working with the Database}
\Gates{test-pubs-examples, test-interp-replace, test-samir-query, test-ndx-maintain}
The programs so far read the ledger.  This chapter adds invoices to it: a
program asks for the invoice, makes sure the number is new, and adds the
record.

\begin{ChapterCovers}
\item The databases a program works on
\item Finding a record with \cmd{SEEK}
\item Adding a record: \cmd{APPEND BLANK} and \cmd{REPLACE}
\item Index files in programs
\item Entering an invoice
\end{ChapterCovers}

\ATEHead{Preparing for This Chapter}
Start InitechBase with a fresh copy of the sample diskette.

\ATEHead{The Databases of a Program}
A program works on the databases opened at the dot prompt before it runs.  The
programs of the system expect \file{INVOICE.DBF} in the selected work area; the
programs that find invoices by number also expect its index \file{INVNO.NDX}.
Each program says so in its comment lines.  A program can select other work
areas with \cmd{SELECT}.

\ATEHead{Index Files in Programs}
The index file named first in \cmd{USE} gives the order in which a program
meets the records, and is the index \cmd{SEEK} searches.  The ledger has two
indexes, on the invoice number and on the customer:
\DotExample{examples/p7-order}
The index files open with a database are kept up to date when a program adds
or changes records.

\ATEHead{Entering an Invoice}
\file{NEWINV.PRG} asks for the invoice number, looks it up, and stops if the
number is already in the ledger.  Otherwise it asks for the rest of the invoice,
adds a blank record and fills it in:
\ProgramListing{NEWINV.PRG}
\begin{Steps}
\item Type \type{USE INVOICE INDEX INVNO}.
\item Type \type{DO NEWINV}.
\item Answer the questions: \type{B1005}, \type{Initrode}, \type{420.00},
\type{04/15/92}.
\item Type \type{LIST INVNO, CUST, AMOUNT, DUEDATE}.
\end{Steps}
\DotExample{examples/p7-newinv}
A number already in the ledger is refused:
\DotExample{examples/p7-dup}

\ATEHead{How the Program Works}
\SideHead{Checking first}The program seeks the new number before it adds
anything.  \cmd{RETURN} inside the \cmd{IF} ends the program at once when the
number is found.

\SideHead{The answers}\cmd{ACCEPT} takes the customer and the due date as
characters, \cmd{INPUT} the amount as a number.

\SideHead{The record}\cmd{APPEND BLANK} adds the record and makes it current;
two \cmd{REPLACE} commands fill in its fields.  \fn{UPPER()} keeps the customer in
capitals and \fn{CTOD()} converts the due date.  The record is written to the
disk at once, and the index is brought up to date.

\ATEHead{Where Do You Go From Here?}
Chapter~8 prints the ledger as a report.

\ATEHead{Summary}
\begin{enumerate}
\item Open the databases and indexes a program needs before it runs.
\item \cmd{SEEK} and \fn{FOUND()} check whether a key is present.
\item \cmd{APPEND BLANK} and \cmd{REPLACE} add a record; the open indexes are
  kept up to date.
\end{enumerate}

% Programming with InitechBase, Chapter 8: Printing (AN-16; annex Sec 8.2
% P12).  Panels replayed (examples/p8-*.dot); sources cmd/report.c (REPORT
% FORM, TO FILE, TO PRINT #126), the report form INVRPT.FRM.
\chapter{Printing}
\Gates{test-pubs-examples, test-samir-report}
The invoice register is the list of every invoice of the ledger with its total,
printed for the files at the end of each month.  The system prints it from a
report form, \file{INVRPT.FRM}, with \cmd{REPORT FORM}.

\begin{ChapterCovers}
\item The report form of the invoice register
\item \cmd{REPORT FORM} in a program
\item Writing a report to a file for the printer
\end{ChapterCovers}

\ATEHead{Preparing for This Chapter}
Start InitechBase with a fresh copy of the sample diskette, and type
\type{USE INVOICE}.

\ATEHead{The Invoice Register}
\file{INVRPT.FRM} gives the report its two-line title and four columns: the
invoice number, the customer, the due date and the amount, which is totalled.
\begin{Steps}
\item Type \type{REPORT FORM INVRPT PLAIN}.
\end{Steps}
\DotExample{examples/p8-register}
\kw{PLAIN} leaves out the page number and the date, so the register can be filed
with the month's papers.

\ATEHead{Reports from Programs}
A program uses \cmd{REPORT FORM} as you do at the dot prompt.  The main menu of
Chapter~11 prints the register when the clerk chooses 3.  A \kw{FOR} condition
prints part of the ledger, and \kw{HEADING} adds a line, such as the month, to
each page.

\ATEHead{Printing on Paper}
\kw{TO FILE} \xp{file} writes the report to a text file, with the extension
\file{.TXT}, as well as displaying it.  Leave InitechBase and copy the file to
the printer with the InitechDOS command \file{COPY}.  \kw{TO PRINT} is refused
with \xmsg{126}{Printer is either not connected or turned off.}

\ATEHead{Where Do You Go From Here?}
Chapter~9 is the first of the month-end programs: the aging report.

\ATEHead{Summary}
\begin{enumerate}
\item \cmd{REPORT FORM} \xp{form} prints a report from a report form file, in a
  program as at the dot prompt.
\item \kw{PLAIN}, \kw{HEADING} and \kw{FOR} change what is printed.
\item \kw{TO FILE} writes the report to a text file for the printer.
\end{enumerate}

% Programming with InitechBase, Chapter 9: The Aging Report (AN-16).  The
% program is Y2KACCT.PRG of the sample diskette, whose body is byte for byte
% the body of harness/diff/dbf_diff/canon/y2k_accounting.prg
% (test-pubs-samir-programs), and its replayed output (examples/p9-aging.dot)
% is that of canon/y2k_accounting.out.  Documented as the shipped program,
% with the date convention of SET CENTURY OFF (CTOD base 1900;
% core/fn_builtins.c fn_ctod_impl; Using InitechBase, SET CENTURY).
\chapter{The Aging Report}
\Gates{test-pubs-examples, test-pubs-samir-programs}
At the end of each month the Collections desk ages the open invoices: for each
invoice it works out how many days have passed since the due date, flags the
invoices more than thirty days past due as overdue, and totals the unpaid
overdue amounts.  \file{Y2KACCT.PRG} is the aging report of the Initech
Accounting System.

\begin{ChapterCovers}
\item The reporting date
\item Date arithmetic
\item The aging report
\item Reading the program
\end{ChapterCovers}

\ATEHead{Preparing for This Chapter}
Start InitechBase with a fresh copy of the sample diskette, and type
\type{USE INVOICE}.

\ATEHead{The Reporting Date}
The report is run as of a reporting date, the memory variable \texttt{ASOF},
which the supervisor sets at the dot prompt before the run.  The report can
therefore be run again for any month-end.  The date is entered, like every date
of the system, in the standard form month, day and year with two digits each,
converted with \fn{CTOD()}.

Under the standard date convention of InitechBase (SET CENTURY OFF), a year
entered with two digits is a year of the nineteen hundreds:
\DotExample{examples/p9-dates}
The due dates of the ledger were entered and stored the same way, on the data
entry screen, so the reporting date and the due dates follow one convention.

\ATEHead{Running the Report}
\begin{Steps}
\item Type \type{ASOF = CTOD('01/31/00')}.
\item Type \type{DO Y2KACCT}.
\end{Steps}
\DotExample{examples/p9-aging}
\SideHead{Age}The age of an invoice is the reporting date minus its due date,
in days.  An invoice whose due date is later than the reporting date has a
negative age: it is not yet due.

\SideHead{Status}An invoice more than thirty days past due is
\texttt{OVERDUE}; every other invoice is \texttt{CURRENT}.

\SideHead{Total}The last line totals the amounts of the overdue invoices that
have not been paid.  For this run every invoice is current, and the total is
zero.

\ATEHead{Reading the Program}
\ProgramListing{Y2KACCT.PRG}
\SideHead{Headings}The first lines print the title and the reporting date,
converted with \fn{DTOC()}.

\SideHead{The loop}For each invoice, \texttt{AGE = ASOF - DUEDATE} subtracts the
dates and gives the age in days.  \fn{IIF()} chooses the status, and \cmd{?}
prints the line, with the age converted by \fn{STR()} to six characters.

\SideHead{The total}\cmd{IF AGE > 30 .AND. .NOT. PAID} adds the amount of each
unpaid overdue invoice to \texttt{OVERTOT}, which is printed after the loop.

\begin{AteNote}
Set \texttt{ASOF} before you run the report.  The report uses the reporting date
throughout, not the date of the computer, so that a month can be closed again
with the same figures.
\end{AteNote}

\ATEHead{Where Do You Go From Here?}
Chapter~10 is the second month-end program: the finance-charge posting.

\ATEHead{Summary}
\begin{enumerate}
\item The reporting date is the memory variable \texttt{ASOF}, set with
  \fn{CTOD()} before the run.
\item Dates entered with two-digit years follow the standard date convention.
\item One date minus another is a number of days; \fn{IIF()} chooses the status.
\end{enumerate}

% Programming with InitechBase, Chapter 10: Finance-Charge Posting (AN-16).
% The program is INTPOST.PRG of the sample diskette, whose body is byte for
% byte the body of harness/diff/dbf_diff/canon/salami.prg
% (test-pubs-samir-programs); its replayed output (examples/p10-intpost.dot)
% is that of canon/salami.out.  Documented as the shipped program.
\chapter{Finance-Charge Posting}
\Gates{test-pubs-examples, test-pubs-samir-programs}
Each month a finance charge is posted to every open invoice.  The charge is
computed to full precision and posted to the customer at the precision of the
billing statement; the difference, the rounding adjustment, is swept into the
suspense account, so that the ledger foots to the cent at close.
\file{INTPOST.PRG} is the finance-charge posting of the Initech Accounting
System.

\begin{ChapterCovers}
\item The posting parameters RATE and SCALE
\item Rounding with \fn{ROUND()}
\item The rounding adjustment and the suspense account
\item Reading the program
\end{ChapterCovers}

\ATEHead{Preparing for This Chapter}
Start InitechBase with a fresh copy of the sample diskette, and type
\type{USE INVOICE}.

\ATEHead{The Posting Parameters}
Two memory variables, set at the dot prompt before the run, govern the posting.
\texttt{RATE} is the monthly finance charge: 0.015, one and one half per cent a
month, is the standard rate.  \texttt{SCALE} is the number of decimal places of
the billing statement: the charge posted to each customer is rounded to that
many places.  Initech's statements are posted in whole dollars, with
\texttt{SCALE} 0.

\ATEHead{Running the Posting}
\begin{Steps}
\item Type \type{RATE = 0.015}.
\item Type \type{SCALE = 0}.
\item Type \type{DO INTPOST}.
\end{Steps}
\DotExample{examples/p10-intpost}
\SideHead{Charge}The finance charge of each open invoice, the amount times the
rate, shown to four decimal places.

\SideHead{Posted}The charge rounded to the statement: \fn{ROUND()} to
\texttt{SCALE} places.

\SideHead{Adjustment}The charge minus the posted amount, positive or negative.

\SideHead{Suspense account}The adjustments of all the postings, added together
in the memory variable \texttt{BOLTON}, the rounding-adjustment suspense
account, for the month-end journal entry.

\ATEHead{Reading the Program}
\ProgramListing{INTPOST.PRG}
The loop visits every invoice and posts only those not paid.  \texttt{CHARGE},
\texttt{POSTED} and \texttt{ADJUST} are computed for each; \texttt{BOLTON} begins at
zero and receives every adjustment.  \fn{STR()} prints the charge and the
adjustment to four places and the posted amount to two.

\begin{AteNote}
Set \texttt{RATE} and \texttt{SCALE} before every run.  The program does not keep
them from one month to the next.
\end{AteNote}

\ATEHead{Where Do You Go From Here?}
Chapter~11 brings the daily tasks of the system together in a menu.

\ATEHead{Summary}
\begin{enumerate}
\item \texttt{RATE} is the monthly rate; \texttt{SCALE} the decimal places of the
  statement.
\item \fn{ROUND()} posts the charge; the remainder is the rounding adjustment.
\item The adjustments are swept into the suspense account \texttt{BOLTON}.
\end{enumerate}

% Programming with InitechBase, Chapter 11: The Main Menu (AN-16).  The
% program is ACCTMAIN.PRG of the sample diskette; panels replayed
% (examples/p11-*.dot).  Sources cmd/flow.c (DO WHILE .T., EXIT, DO CASE),
% cmd/proc.c (WAIT, procedures), cmd/report.c.
\chapter{The Main Menu}
\Gates{test-pubs-examples, test-interp-flow, test-interp-proc}
The accounts receivable clerk uses the system every day to look at the ledger
and print the register.  The main menu offers these tasks by number and carries
out the one chosen, until the clerk chooses to leave.

\begin{ChapterCovers}
\item A menu loop
\item Taking the choice with \cmd{WAIT}
\item Carrying it out with \cmd{DO CASE}
\item Procedures of the menu
\end{ChapterCovers}

\ATEHead{Preparing for This Chapter}
Start InitechBase with a fresh copy of the sample diskette, and type
\type{USE INVOICE INDEX INVNO}.

\ATEHead{Using the Menu}
\begin{Steps}
\item Type \type{DO ACCTMAIN}.
\ATEResponse{The menu appears and waits for a choice.}
\item Press \type{2}.
\ATEResponse{The open invoices are listed, and the menu appears again.}
\item Press \type{9}.
\ATEResponse{The menu asks for one of its choices.}
\item Press \type{0}.
\ATEResponse{The dot prompt returns.}
\end{Steps}
\DotExample{examples/p11-menu}
Choice 1 lists the whole ledger:
\DotExample{examples/p11-ledger}

\ATEHead{Reading the Program}
\ProgramListing{ACCTMAIN.PRG}
\SideHead{The loop}\cmd{DO WHILE .T.} repeats for as long as the clerk works;
the condition is always true.  The loop ends only by \cmd{EXIT}, in the
\cmd{CASE} of choice 0.

\SideHead{The menu}\cmd{?} prints an empty line, the title and the choices.  A
second empty \cmd{?} begins a new line for the question.

\SideHead{The choice}\cmd{WAIT} takes one key into \texttt{MCHOICE}.

\SideHead{The task}\cmd{DO CASE} compares the key with each choice and carries out
its task: \cmd{DO LEDGER}, \cmd{DO OPENLIST}, or \cmd{REPORT FORM}.  Any other key
reaches \cmd{OTHERWISE}, which asks again.

\SideHead{The procedures}\texttt{LEDGER} and \texttt{OPENLIST} follow the main
program in the same file.  Each begins at the first record, because the task
before it may have left the record pointer at the end of the file.

\ATEHead{Where Do You Go From Here?}
Chapter~12 describes what happens when a command of a program cannot be
carried out.

\ATEHead{Summary}
\begin{enumerate}
\item \cmd{DO WHILE .T.} \ldots{} \cmd{ENDDO} with \cmd{EXIT} makes a menu loop.
\item \cmd{WAIT} \ldots{} \kw{TO} takes the choice; \cmd{DO CASE} carries it out.
\item The tasks are procedures of the same file.
\end{enumerate}

% Programming with InitechBase, Chapter 12: Handling Errors (AN-16).
% Panels replayed (examples/p12-*.dot); sources samir_main.c (the message
% catalogue, "n  text"), cmd/proc.c (ON ERROR, proc_run), core/fn_builtins.c
% ERROR().  The program BADTOTAL.PRG of the sample diskette.
\chapter{Handling Errors}
\Gates{test-pubs-examples, test-interp-proc, test-samir-repl}
A command that cannot be carried out stops the program.  This chapter
describes the messages InitechBase displays, how to find the command at fault,
and how a program can record an error with \cmd{ON ERROR}.

\begin{ChapterCovers}
\item Error messages and their numbers
\item What happens to a program
\item \cmd{ON ERROR} and \fn{ERROR()}
\end{ChapterCovers}

\ATEHead{Preparing for This Chapter}
Start InitechBase with a fresh copy of the sample diskette.

\ATEHead{Error Messages}
When a command cannot be carried out, InitechBase displays a message: its
number, two blanks and its text.  Three common errors at the dot prompt:
\DotExample{examples/p12-errors}
\xmsg{9}{Data type mismatch.} means that an operator was given values of types
it cannot combine; \xmsg{12}{Variable not found.} that a name is neither a field
of the database in USE nor a memory variable; and the third that a function was
given an argument out of its range.  Chapter~7 of \doctitle{Using InitechBase}
lists every message with its cause and what to do.

\ATEHead{What Happens to a Program}
When a command of a program fails, the program stops at that command.
InitechBase displays the message and the dot prompt.  The memory variables the
program created remain, and the record pointer stays where it was; both help
you find the command at fault (Chapter~13).

\file{BADTOTAL.PRG} totals the ledger, but its programmer misspelled the field
\texttt{AMOUNT}:
\ProgramListing{BADTOTAL.PRG}

\ATEHead{ON ERROR}
\cmd{ON ERROR} names a command to be carried out whenever an error occurs.  While
the command runs, \fn{ERROR()} gives the number of the error.  A program uses it
to record the error before the message is displayed:
\begin{Steps}
\item Type \type{USE INVOICE}.
\item Type \type{ON ERROR STORE ERROR() TO MERR}.
\item Type \type{DO BADTOTAL}.
\item Type \type{? MERR, RECNO()}, then \type{ON ERROR} to remove the trap.
\end{Steps}
\DotExample{examples/p12-onerror}
The program stopped on its first record, with error 12.

\ATEHead{Where Do You Go From Here?}
Chapter~13 shows how to find and correct the mistake.

\ATEHead{Summary}
\begin{enumerate}
\item An error message gives a number and a text; Using InitechBase, Chapter~7,
  explains each.
\item A program stops at the command that fails; its memory variables remain.
\item \cmd{ON ERROR} \xp{command} runs a command when an error occurs;
  \fn{ERROR()} gives its number.
\end{enumerate}

% Programming with InitechBase, Chapter 13: Testing and Debugging (AN-16;
% annex Sec 8.2 P15).  Panels replayed (examples/p13-*.dot); the program
% BADTOTAL.PRG; the memory variables of a main program remain after it
% (cmd/flow.c level 0, observed).
\chapter{Testing and Debugging}
\Gates{test-pubs-examples, test-interp-proc}
A program that runs is not yet a program that is right.  This chapter gives the
habits that find mistakes early, and shows how to find the command at fault
when a program stops.

\begin{ChapterCovers}
\item Testing on copies
\item Looking at a program after it ends
\item Finding the command at fault
\item Testing a part of a program at the dot prompt
\end{ChapterCovers}

\ATEHead{Preparing for This Chapter}
Start InitechBase with a fresh copy of the sample diskette.

\ATEHead{Testing on Copies}
Test a program on a copy of the databases, never on the files of the business.
A program that changes records -- \file{NEWINV.PRG}, for one -- writes to the disk
at once.  Keep a copy of the files as they were before the test, so that you can
run it again from the same start.

\ATEHead{Looking After the Run}
The memory variables a program creates remain after it ends, and so does the
record pointer.  Look at them with \cmd{?}:
\DotExample{examples/p13-check}
\fn{TYPE()} tells you whether a memory variable has the type you meant; a total
that turns out to be of type C was given a character value somewhere.

\ATEHead{Finding the Command at Fault}
\file{BADTOTAL.PRG}, listed in Chapter~12, stops with a message:
\DotExample{examples/p13-bad}
The message says that a name was not found.  \texttt{MTOTAL} is still 0 and the
record pointer is on record 1: the program stopped inside the loop, on the
first record, before anything was added.  The only name used there that is not
a memory variable is \texttt{AMOUNTS}: the field is \texttt{AMOUNT}.  Correct the
line with your editor and run the program again.

\ATEHead{Testing a Part at the Dot Prompt}
Every command of a program can be typed at the dot prompt, except the
structures \cmd{IF}, \cmd{DO CASE} and \cmd{DO WHILE}.  To test a calculation,
open the database, move to a record, store the memory variables the
calculation needs, and type the calculation with \cmd{?}.  When it gives the
right result, copy it into the program.

\begin{Tip}
Add a \cmd{?} to a program to display a memory variable at the point where you
suspect the mistake, and remove it when the program is right.
\end{Tip}

\ATEHead{Where Do You Go From Here?}
Chapter~14 works with two databases at once.

\ATEHead{Summary}
\begin{enumerate}
\item Test on copies of the databases.
\item Memory variables and the record pointer remain after a program; look at
  them with \cmd{?} and \fn{TYPE()}.
\item The message, the record number and the memory variables lead you to the
  command at fault.
\item Try calculations at the dot prompt before you put them in a program.
\end{enumerate}

% Programming with InitechBase, Chapter 14: Working with Two Databases
% (AN-16; annex Sec 8.2 P10).  Panels replayed (examples/p14-*.dot);
% sources cmd/flow.c SELECT, cmd/workarea.c (ten areas, each its own
% pointer), cmd/query.c SEEK; a field of another area is reached through a
% memory variable (the shipped evaluator resolves names in the selected
% area).
\chapter{Working with Two Databases}
\Gates{test-pubs-examples, test-interp-use, test-samir-query}
The departments of Initech are kept in their own database, \file{DEPTS.DBF},
with the name and the floor of each; the staff roster gives only the code.  A
program that prints an employee with the name of the department looks the code
up in the second database.  This chapter shows the way InitechBase programs do
it: two work areas, a memory variable, and \cmd{SEEK}.

\begin{ChapterCovers}
\item Two databases in two work areas
\item Carrying a value from one to the other
\item Looking up a code with \cmd{SEEK}
\item Counting the records of another database
\end{ChapterCovers}

\ATEHead{Preparing for This Chapter}
Start InitechBase with a fresh copy of the sample diskette.

\ATEHead{Two Work Areas}
Each database is opened in its own work area, at the dot prompt, before the
program runs.  The program selects the work area of the database it wants to
read; each work area keeps its own record pointer, so moving in one does not
move the other.

A field name always refers to the database of the selected work area.  To use
a value of the staff roster while the departments are selected, the program
stores it in a memory variable first.

\ATEHead{Looking Up the Department}
\file{STAFFDEP.PRG} lists every employee with the name of the department:
\ProgramListing{STAFFDEP.PRG}
\begin{Steps}
\item Type \type{SELECT 2}, then \type{USE DEPTS INDEX DEPTS}.
\item Type \type{SELECT 1}, then \type{USE STAFF}.
\item Type \type{DO STAFFDEP}.
\end{Steps}
\DotExample{examples/p14-staffdep}
For each employee the program keeps the code in \texttt{MDEPT}, selects the
departments, seeks the code in their index and takes the name; then it selects
the staff again before \cmd{SKIP}.  A code with no department would print the
words in parentheses instead of a name.

\ATEHead{Counting Through an Index}
\file{HEADCNT.PRG} works the other way round: for each department it counts the
employees, through the index of the staff on the department:
\ProgramListing{HEADCNT.PRG}
\DotExample{examples/p14-headcnt}
The departments are read in their order on the disk; for each, the inner loop
seeks the first employee of the department and counts while the department
stays the same.

\ATEHead{Where Do You Go From Here?}
Chapter~15 prints a long listing in pages.

\ATEHead{Summary}
\begin{enumerate}
\item Open each database in its own work area before the program runs.
\item \cmd{SELECT} chooses the database whose fields a name refers to.
\item Carry a value from one database to another in a memory variable.
\item \cmd{SEEK} looks the value up in the index of the other database.
\end{enumerate}

% Programming with InitechBase, Chapter 15: Reports from Programs (AN-16;
% annex Sec 8.2 P12).  Panel replayed (examples/p15-pagerpt.dot); the
% program PAGERPT.PRG of the sample diskette.
\chapter{Reports from Programs}
\Gates{test-pubs-examples, test-interp-flow, test-interp-proc}
A report form lays out a report in columns.  When a report needs a layout of
its own -- lines built from several fields, a heading of its own, a count of
lines -- a program prints it with \cmd{?}.  This chapter prints the staff
roster in numbered pages.

\begin{ChapterCovers}
\item Counting lines
\item A heading on every page
\item Page numbers
\end{ChapterCovers}

\ATEHead{Preparing for This Chapter}
Start InitechBase with a fresh copy of the sample diskette, and type
\type{USE STAFF}.

\ATEHead{Pages of Lines}
A page holds a fixed number of lines.  The program counts the lines it has
printed on the page in a memory variable; when the page is full, it prints the
heading of a new page and begins the count again.  Starting the count at the
full value makes the program print the heading before the first line, too.

\ATEHead{The Roster in Pages}
\file{PAGERPT.PRG} prints six employees to a page, so that the example fits this
book; for a printed page of 55 lines, set \texttt{MLINES} to 55.
\ProgramListing{PAGERPT.PRG}
\begin{Steps}
\item Type \type{DO PAGERPT}.
\end{Steps}
\DotExample{examples/p15-pagerpt}
\SideHead{MLINE}counts the lines of the current page; \texttt{MPAGE} the pages.

\SideHead{HEADING}is the procedure that begins a page: an empty line, the title
with the page number, the column headings, and the count set back to zero.

\SideHead{The end}After the loop, a closing line tells the reader that the
roster is complete.

\ATEHead{Where Do You Go From Here?}
Chapter~16 calculates with dates.

\ATEHead{Summary}
\begin{enumerate}
\item Count the lines printed on a page in a memory variable.
\item Print the heading when the count reaches the length of the page.
\item Keep the page number in a memory variable and print it in the heading.
\end{enumerate}

% Programming with InitechBase, Chapter 16: Working with Dates (AN-16).
% Panels replayed (examples/p16-*.dot); sources core/eval.c (date - date,
% date + n), core/fn_builtins.c (CTOD, DOW, CDOW, DAY), the program
% SERVICE.PRG.
\chapter{Working with Dates}
\Gates{test-pubs-examples, test-xbase-fn-b, test-xbase-eval}
Personnel records are full of dates: of hire, of review, of birth.  This
chapter computes the years of service of every employee, as of a date the
program is given.

\begin{ChapterCovers}
\item Date arithmetic
\item The day of the week
\item Years of service
\end{ChapterCovers}

\ATEHead{Preparing for This Chapter}
Start InitechBase with a fresh copy of the sample diskette, and type
\type{USE STAFF}.

\ATEHead{Date Arithmetic}
One date minus another is the number of days between them; a date plus or
minus a number of days is a date.  The calendar is followed exactly, leap years
included:
\DotExample{examples/p16-dates}
\fn{DOW()} numbers the days of the week from 1, Sunday, to 7, Saturday;
\fn{CDOW()} names them.

\ATEHead{Years of Service}
\file{SERVICE.PRG} computes, for each employee, the complete years of service
as of the date in \texttt{MASOF}, and counts the staff in three classes:
\ProgramListing{SERVICE.PRG}
\begin{Steps}
\item Type \type{MASOF = CTOD('03/01/92')}.
\item Type \type{DO SERVICE}.
\end{Steps}
\DotExample{examples/p16-service}
The days of service divided by 365, with the fraction dropped by \fn{INT()},
give the years.  Like the aging report of Chapter~9, the program takes its date
from a memory variable, so it can be run as of any date.

\ATEHead{Where Do You Go From Here?}
Chapter~17 works with the characters of names.

\ATEHead{Summary}
\begin{enumerate}
\item Date minus date gives days; date plus days gives a date.
\item \fn{DOW()} and \fn{CDOW()} give the day of the week.
\item \fn{INT()} of the days divided by 365 gives complete years.
\end{enumerate}

% Programming with InitechBase, Chapter 17: Working with Characters (AN-16).
% Panels replayed (examples/p17-*.dot); sources core/fn_builtins.c (LEFT,
% SUBSTR, LOWER, UPPER, TRIM, AT, STUFF), core/eval.c ($), the program
% NAMEFMT.PRG.
\chapter{Working with Characters}
\Gates{test-pubs-examples, test-xbase-fn-a, test-xbase-fn-c}
The staff roster keeps every name in capitals, as the clerks key them.  A
letter to an employee prints the name with only its first letter in capitals.
This chapter builds the salutation of a letter from the fields of the roster.

\begin{ChapterCovers}
\item Taking strings apart: \fn{LEFT()} and \fn{SUBSTR()}
\item Changing case: \fn{UPPER()} and \fn{LOWER()}
\item Searching: \texttt{\$} and \fn{AT()}
\item Replacing part of a string: \fn{STUFF()}
\end{ChapterCovers}

\ATEHead{Preparing for This Chapter}
Start InitechBase with a fresh copy of the sample diskette, and type
\type{USE STAFF}.

\ATEHead{The Functions}
\texttt{\$} is true when one string occurs in another; \fn{AT()} gives the place
where it begins; \fn{STUFF()} replaces part of a string with another:
\DotExample{examples/p17-strings}
The first letter of a name, in capitals, joined to the rest in small letters,
gives the name as a letter prints it.

\ATEHead{Salutations}
\file{NAMEFMT.PRG} prints the salutation of a letter to each employee:
\ProgramListing{NAMEFMT.PRG}
\begin{Steps}
\item Type \type{DO NAMEFMT}.
\end{Steps}
\DotExample{examples/p17-namefmt}
\fn{TRIM()} removes the blanks at the end of each field, so that the first and
the last name are separated by one blank and the comma follows the name.

\ATEHead{Where Do You Go From Here?}
Chapter~18 changes many records in one program.

\ATEHead{Summary}
\begin{enumerate}
\item \fn{LEFT()} and \fn{SUBSTR()} take strings apart; \texttt{+} joins them.
\item \fn{UPPER()} and \fn{LOWER()} change case; \fn{TRIM()} removes trailing
  blanks.
\item \texttt{\$}, \fn{AT()} and \fn{STUFF()} search and replace.
\end{enumerate}

% Programming with InitechBase, Chapter 18: Changing Many Records (AN-16).
% Panels replayed (examples/p18-*.dot); sources cmd/mutate.c (REPLACE on
% the current record, written at once), the programs RAISEDPT.PRG and
% PAYTOTAL.PRG.
\chapter{Changing Many Records}
\Gates{test-pubs-examples, test-interp-replace}
Once a year Initech raises its salaries, by a percentage that depends on the
department.  A single \cmd{REPLACE} with \kw{FOR} gives every record the same
raise; a program gives each department its own.

\begin{ChapterCovers}
\item Changing records one at a time in a loop
\item A value that depends on the record
\item Showing every change
\item Checking the result
\end{ChapterCovers}

\ATEHead{Preparing for This Chapter}
Start InitechBase with a fresh copy of the sample diskette, and type
\type{USE STAFF}.  The program of this chapter changes the roster: copy the
diskette again afterwards.

\ATEHead{Raises by Department}
\file{RAISEDPT.PRG} chooses the percentage of each employee with \cmd{DO CASE},
keeps the old salary, replaces it, and prints both:
\ProgramListing{RAISEDPT.PRG}
\begin{Steps}
\item Type \type{DO RAISEDPT}.
\end{Steps}
\DotExample{examples/p18-raisedpt}
\fn{ROUND()} keeps every new salary to the cent.  The total added to the
payroll is the sum of the differences.

\ATEHead{Checking the Result}
Run a report of the totals after the change.  \file{PAYTOTAL.PRG}, of Volume
One, totals the payroll of the exempt and the non-exempt staff:
\DotExample{examples/p18-after}
\begin{AteWarning}
A program that replaces fields changes the database at once and cannot undo the
change.  Run it on a copy, and keep the copy until the result has been checked.
\end{AteWarning}

\ATEHead{Where Do You Go From Here?}
Chapter~19 reads the memo fields of the roster.

\ATEHead{Summary}
\begin{enumerate}
\item A loop with \cmd{REPLACE} gives each record its own new value.
\item Keep the old value in a memory variable to show or check the change.
\item Check the totals after the change, and work on a copy.
\end{enumerate}

% Programming with InitechBase, Chapter 19: Putting It Together (AN-16;
% annex Sec 8.2 P14).  The programs and files of the Initech Accounting
% System as the sample diskette holds them (pubs/examples/samir/disk), and
% the daily and month-end procedures with the commands of Chapters 1-13.
\chapter{Putting It Together}
\Gates{test-pubs-build, test-pubs-samir-disk}
This chapter describes the Initech Accounting System as a whole: its files, its
programs, and the procedures by which the clerk and the supervisor use it each
day and at the end of each month.

\IsThisGoodForTheCompany

\ATEHead{The Files}
\noindent\hspace*{-\imHang}\begin{tabular}{@{}p{30mm}p{104mm}@{}}
\hline
\textsf{\textbf{File}} & \textsf{\textbf{Contents}}\\
\hline
\file{INVOICE.DBF} & The accounts receivable ledger: \texttt{INVNO} C 5,
  \texttt{CUST} C 10, \texttt{AMOUNT} N 10.2, \texttt{DUEDATE} D, \texttt{PAID} L\\
\file{INVNO.NDX} & Index of the ledger on the invoice number\\
\file{INVCUST.NDX} & Index of the ledger on the customer\\
\file{INVRPT.FRM} & Report form of the invoice register\\
\hline
\end{tabular}

\ATEHead{The Programs}
\noindent\hspace*{-\imHang}\begin{tabular}{@{}p{30mm}p{60mm}p{40mm}@{}}
\hline
\textsf{\textbf{Program}} & \textsf{\textbf{Task}} & \textsf{\textbf{Chapter}}\\
\hline
\file{ACCTMAIN.PRG} & The main menu of the clerk & 11\\
\file{OPENINV.PRG} & List of the open invoices & 1\\
\file{INVTOTAL.PRG} & Totals of the ledger & 2, 3\\
\file{LATEFEE.PRG} & Late fee of one invoice & 4\\
\file{INVFMT.PRG} & Formatted register & 5\\
\file{FINDINV.PRG} & Finding an invoice & 6\\
\file{NEWINV.PRG} & Entering an invoice & 7\\
\file{Y2KACCT.PRG} & Month-end aging report & 9\\
\file{INTPOST.PRG} & Month-end finance-charge posting & 10\\
\hline
\end{tabular}

\ATEHead{Each Day}
\begin{Steps}
\item Start InitechBase from the disk that holds the system.
\item Type \type{USE INVOICE INDEX INVNO}.
\item Type \type{DO ACCTMAIN}, and work from the menu.
\item Enter new invoices with \type{DO NEWINV} at the dot prompt.
\item Leave InitechBase with \type{QUIT}.
\end{Steps}

\ATEHead{At the End of the Month}
\begin{Steps}[resume]
\item Copy the files of the system to a diskette kept for the month.
\item Type \type{USE INVOICE}.
\item Set the reporting date, for example \type{ASOF = CTOD('01/31/00')}, and
  type \type{DO Y2KACCT}.  File the aging report with Collections.
\item Type \type{RATE = 0.015} and \type{SCALE = 0}, then \type{DO INTPOST}.
  Record the suspense account total in the month-end journal.
\item Type \type{REPORT FORM INVRPT PLAIN TO FILE INVRPT}, leave InitechBase, and
  print \file{INVRPT.TXT} for the files.
\end{Steps}

\ATEHead{Where Do You Go From Here?}
You have built and used a complete InitechBase application.  For the exact form
of any command or function, and for the meaning of every message, turn to
\doctitle{Using InitechBase} in Volume One.

\ATEHead{Summary}
\begin{enumerate}
\item The ledger \file{INVOICE.DBF}, its indexes and its report form are the
  files of the system.
\item The clerk works from the main menu each day.
\item At month-end the supervisor copies the files, runs the aging report and the
  finance-charge posting, and prints the register.
\end{enumerate}

\appendix
% Programming with InitechBase, Appendix A: Commands at the Dot Prompt and
% in Programs.  Sources: samir_main.c samir_repl (USE, CLOSE, QUIT, EXIT and
% DO of a program file are carried out by the dot prompt itself; observed
% #16 when a program uses them), cmd/flow.c (IF, DO CASE, DO WHILE and their
% parts need a program: at the dot prompt they are refused), cmd/proc.c.
\chapter{Commands at the Dot Prompt and in Programs}
\Gates{test-samir-repl, test-interp-flow, test-interp-proc}
Most commands may be typed at the dot prompt and written in programs alike.
A few belong to one or the other.

\ATEHead{Only at the Dot Prompt}
\noindent\hspace*{-\imHang}\begin{tabular}{@{}p{38mm}p{96mm}@{}}
\hline
\textsf{\textbf{Command}} & \textsf{\textbf{In a program}}\\
\hline
\cmd{USE} & Open the databases at the dot prompt before the program runs.\\
\cmd{CLOSE} & Close the databases at the dot prompt after the program.\\
\cmd{QUIT}, \cmd{EXIT} & Leave InitechBase from the dot prompt.  In a program,
  \cmd{EXIT} leaves a \cmd{DO WHILE} loop.\\
\cmd{DO} \xp{program file} & A program calls the procedures of its own file;
  program files are run from the dot prompt.\\
\hline
\end{tabular}

\ATEHead{Only in Programs}
\noindent\hspace*{-\imHang}\begin{tabular}{@{}p{38mm}p{96mm}@{}}
\hline
\textsf{\textbf{Command}} & \textsf{\textbf{Use}}\\
\hline
\cmd{IF}, \cmd{ELSE}, \cmd{ENDIF} & Decisions\\
\cmd{DO CASE}, \cmd{CASE}, \cmd{OTHERWISE}, \cmd{ENDCASE} & Choices\\
\cmd{DO WHILE}, \cmd{LOOP}, \cmd{EXIT}, \cmd{ENDDO} & Loops\\
\cmd{PROCEDURE} & Begins a procedure of a program file\\
\cmd{PARAMETERS} & Receives the values of \cmd{DO} \ldots{} \kw{WITH}\\
\hline
\end{tabular}

\ATEHead{Both}
Every other command of \doctitle{Using InitechBase}, Chapter~5, may be typed at
the dot prompt and written in a program: \cmd{?}, \cmd{??}, \cmd{ACCEPT},
\cmd{APPEND BLANK}, \cmd{CONTINUE}, \cmd{DELETE}, \cmd{DISPLAY}, \cmd{DO} of a
procedure, \cmd{FIND}, \cmd{GO}, \cmd{INPUT}, \cmd{LIST}, \cmd{LOCATE},
\cmd{NOTE}, \cmd{ON ERROR}, \cmd{PACK}, \cmd{PRIVATE}, \cmd{PUBLIC},
\cmd{RECALL}, \cmd{RELEASE}, \cmd{REPLACE}, \cmd{REPORT FORM}, \cmd{RETURN},
\cmd{SEEK}, \cmd{SELECT}, \cmd{SKIP}, \cmd{STORE}, \cmd{WAIT}, \cmd{ZAP} and
the SET commands.

% Programming with InitechBase, Appendix B: The Sample Programs.  The
% programs of the Sample Programs and Utilities diskette
% (pubs/examples/samir/disk/MANIFEST), with the database each expects and
% the chapter that describes it.
\chapter{The Sample Programs}
\Gates{test-pubs-samir-disk}
\noindent\hspace*{-\imHang}\begin{tabular}{@{}p{26mm}p{52mm}p{34mm}p{8mm}@{}}
\hline
\textsf{\textbf{Program}} & \textsf{\textbf{Task}} & \textsf{\textbf{Opened first}} &
  \textsf{\textbf{Ch.}}\\
\hline
\file{ACCTMAIN.PRG} & Main menu of the clerk & \file{INVOICE}, \file{INVNO} & 11\\
\file{AMTCHECK.PRG} & Asking until the answer is right & -- & 6\\
\file{BADTOTAL.PRG} & Ledger total with a misspelled field & \file{INVOICE} & 12, 13\\
\file{CLEANUP.PRG} & Removing paid invoices & \file{INVOICE} & 7\\
\file{CUSTTOT.PRG} & Totals by customer & \file{INVOICE}, \file{INVCUST} & 7\\
\file{DUECLASS.PRG} & Invoices by size & \file{INVOICE} & 2\\
\file{ERRDEMO.PRG} & An error procedure & \file{INVOICE} & 12\\
\file{FINDINV.PRG} & Finding an invoice & \file{INVOICE}, \file{INVNO} & 6\\
\file{FIRSTBIG.PRG} & First large open invoice & \file{INVOICE} & 2\\
\file{INTPOST.PRG} & Month-end finance-charge posting & \file{INVOICE} & 10\\
\file{INVFMT.PRG} & Formatted register & \file{INVOICE} & 5\\
\file{INVTOTAL.PRG} & Totals of the ledger & \file{INVOICE} & 2, 3\\
\file{LATEFEE.PRG} & Late fee of one invoice & \file{INVOICE} & 4\\
\file{MARKPAID.PRG} & Recording a payment & \file{INVOICE}, \file{INVNO} & 7\\
\file{NEWINV.PRG} & Entering an invoice & \file{INVOICE}, \file{INVNO} & 7\\
\file{OPENINV.PRG} & List of the open invoices & \file{INVOICE} & 1, 2\\
\file{STATEMNT.PRG} & Customer statement & \file{INVOICE}, \file{INVCUST} & 4\\
\file{TITLES.PRG} & Centred titles & -- & 5\\
\file{Y2KACCT.PRG} & Month-end aging report & \file{INVOICE} & 9\\
\hline
\end{tabular}

The programs \file{STAFFLST.PRG}, \file{PAYTOTAL.PRG}, \file{DEPTNAME.PRG},
\file{RAISE.PRG} and \file{TALLY.PRG} work on \file{STAFF.DBF}; they are
described in Volume One.

% Programming with InitechBase, Appendix C: Programming Conventions of the
% Initech Accounting System -- the rules every program of the sample
% diskette follows (pubs/examples/samir/disk/*.PRG).
\chapter{Programming Conventions}
\Gates{test-pubs-build}
The programs of the Initech Accounting System are written to the same rules,
so that any programmer of the Finance and Accounting section can read and
maintain any of them.  Follow the same rules in your own programs.

\SideHead{Comment header}Every program begins with comment lines that give its
file name and title, the section that owns it, the database and index files it
expects to find open, and in a sentence or two what it does and how it is run.

\SideHead{Names}Memory variables begin with M (\texttt{MTOTAL}); parameters with P
(\texttt{PAMOUNT}).  No memory variable has the name of a field.

\SideHead{Structure}Each structure is indented by two blanks.  A program that is
more than a screen long is divided into procedures with names that say what
they do.

\SideHead{Lines}One command to a line, and no line longer than the screen.  A
long expression is built in steps, in memory variables.

\SideHead{Checking first}A program checks before it changes: that the record
exists, that it may be changed, and, for a change that cannot be undone, that
the operator confirms.

\SideHead{Questions}Character answers are taken with \cmd{ACCEPT} and converted to
capitals; amounts with \cmd{INPUT}, checked with \fn{TYPE()}; choices with
\cmd{WAIT}.

\SideHead{Dates}Dates are entered in the standard form month, day and year,
two digits each, and converted with \fn{CTOD()}.  A report that is run as of a
date takes the date from a memory variable set before the run.

\SideHead{The end}Every program and procedure ends with \cmd{RETURN}.

\SectionNotes{Programming Notes}

\EndSections
\IndexKey{\IndexKeyLine{G}{Getting Started (Volume One)}\IndexKeyLine{L}{Learning InitechBase (Volume One)}%
  \IndexKeyLine{U}{Using InitechBase (Volume One)}\IndexKeyLine{P}{Programming with InitechBase}}
\printindex
\FinishBook[3]
\end{document}
